\documentclass[uplatex,a4paper,dvipdfmx,aspectratio=169,11pt]{beamer}

\usepackage{pxjahyper}
\usepackage{amsmath,amsfonts,amssymb,amsthm,bm,ascmac}
\usepackage{ascmac}
\usepackage{physics}
\usepackage{tikz}
\usepackage{gnuplot-lua-tikz}
\usetikzlibrary{intersections,calc,arrows.meta,cd,automata,positioning}
\usepackage{circuitikz}
\usepackage{here}
\usepackage{siunitx}
\usepackage{multicol,multirow}
\usepackage{systeme}
\usepackage[version=3]{mhchem}
\usepackage{chemfig}
\usepackage{url}
\usepackage{braket}
\usepackage{enumerate}
\usepackage{mathrsfs}
\usepackage{otf}
\usepackage{ulem}
\usepackage{stmaryrd}
\usepackage{listings}
\usepackage{bussproofs}
\usepackage{mathtools}
\usepackage{cite}
\usepackage{docmute}
\usepackage{xcolor}
\usepackage{wasysym}
\usepackage{pxrubrica}
\DeclareMathOperator{\Sinarc}{\mathrm{Sin}^{-1}}
\DeclareMathOperator{\Cosarc}{\mathrm{Cos}^{-1}}
\DeclareMathOperator{\Tanarc}{\mathrm{Tan}^{-1}}
% \DeclareMathOperator{\Real}{\bm{R}}
% \DeclareMathOperator{\Complex}{\bm{C}}
% \DeclareMathOperator{\Rational}{\bm{Q}}
% \DeclareMathOperator{\Natural}{\bm{N}}
% \DeclareMathOperator{\Integer}{\bm{Z}}
\DeclareBoldMathCommand{\Real}{R}
\DeclareBoldMathCommand{\Complex}{C}
\DeclareBoldMathCommand{\Rational}{Q}
\DeclareBoldMathCommand{\Natural}{N}
\DeclareBoldMathCommand{\Integer}{Z}
\DeclareMathOperator{\Ker}{\mathrm{Ker}}
\DeclareMathOperator{\diffset}{\backslash}
\DeclareMathOperator{\define}{\overset{\text{def}}{\Longleftrightarrow}}
\newcommand{\typed}{\!:\!}
\newcommand{\Var} {\mathord{\mathbf{Var}}}
\newcommand{\TVar}{\mathord{\mathbf{TVar}}}
\newcommand{\Type}{\mathord{\mathbf{Type}}}
\newcommand{\Term}{\mathord{\mathbf{\Lambda}}}
\newcommand{\Asn}{\mathord{\mathbf{Asn}}}
\newcommand{\Cont}{\mathord{\mathbf{Cont}}}
\newcommand{\Church}{\mathord{\mathbf{M}}}
\newcommand{\Sub}  {\mathord{\mathrm{Sub}}}
\newcommand{\FV}  {\mathord{\mathrm{FV}}}
\newcommand{\BV}  {\mathord{\mathrm{BV}}}
\newcommand{\vars}{\mathcal{V}}
\newcommand{\tvars}{\mathbb{V}}
\newcommand{\terms}{\mathcal{T}}
\newcommand{\types}{\mathbb{T}}
\newcommand{\contexts}{\mathfrak{G}}
\newcommand{\ID}  {(\mathrm{ID})}
\newcommand{\WEAK}{(\mathrm{WEAK})}
\newcommand{\EXC}{(\mathrm{EXC})}
\newcommand{\VAR} {(\mathrm{VAR})}
\newcommand{\APP} {(\mathrm{APP})}
\newcommand{\BND} {(\mathrm{BND})}
\newcommand{\INST}{(\mathrm{INST})}
\newcommand{\GEN} {(\mathrm{GEN})}
\newcommand{\tbeta}{\to_\beta}
\newcommand{\ttbeta}{\twoheadrightarrow_\beta}
\newcommand{\dom}{\mathord{\mathrm{dom}}}
\newcommand{\Typ}{\mathord{\mathrm{Typ}}}
\newcommand{\SN}{\mathsf{SN}}
\newcommand{\SAT}{\mathsf{SAT}}
\newcommand{\Bool}{\mathbb{B}}
\newcommand{\TCP}{\mathbf{TCP}}
\newcommand{\SUP}{\mathbf{SUP}}
\newcommand{\Ltt}{\mathtt{L}}
\newcommand{\Rtt}{\mathtt{R}}
\newcommand{\Stt}{\mathtt{S}}
% \newcommand{\init}{\mathrm{init}}
\newcommand{\init}{0}
\newcommand{\acc}{\mathrm{acc}}
\newcommand{\rej}{\mathrm{rej}}
\newcommand{\bak}{\mathrm{bak}}
\newcommand{\textspace}{\text{\textvisiblespace}\,}
\newcommand{\terminate}{\!\downarrow}
\newcommand{\nterminate}{\!\uparrow}
\newcommand{\problems}{\mathcal{P}_\Sigma}
\newcommand{\classified}{\in}
\newcommand{\Halt}{\mathtt{Halt}}
\newcommand{\ord}{\mathord{\mathrm{ord}}}
\newcommand{\OrdTypes}{\mathord{\mathrm{OrdTypes}}}
\newcommand{\power}[1]{\mathord{\mathfrak{P}}\qty(#1)}
\newcommand{\nbhs}[1]{\mathord{\mathfrak{N}}\qty(#1)}
\newcommand{\Opens}{\mathscr{O}}
\newcommand{\Closeds}{\mathfrak{A}}
\newcommand{\inv}{{\mathrm{-1}}}

\DeclarePairedDelimiter{\interpret}{\llbracket}{\rrbracket}
\DeclarePairedDelimiter{\encode}{\langle}{\rangle}

\DeclareMathOperator{\betared}{\rightarrow_\beta}
\DeclareMathOperator{\lbetared}{\longrightarrow_\beta}
\DeclareMathOperator{\tbetared}{\twoheadrightarrow_\beta}
\newcommand{\nat}{\mathtt{nat}}
\newcommand{\snat}{\nat_\sigma}
\newcommand{\CSet}{{\mathbf{Set}}}
\DeclareMathOperator{\Hom}{{\mathrm{Hom}}}
\newcommand{\lamto}{{\lambda\!\!\to}}
\newcommand{\lamP}{\lambda\mathrm{P}}
\newcommand{\lamT}{\lambda 2}
\newcommand{\lamPT}{\lambda\mathrm{P}2}
\newcommand{\lamo}{\lambda\underline{\omega}}
\newcommand{\lamPo}{\lambda\mathrm{P}\underline{\omega}}
\newcommand{\lamO}{\lambda\omega}
\newcommand{\lamC}{\lambda\mathrm{C}}

\DeclareMathOperator{\pterm}{\mathbf{Pt}}
\newcommand{\PROP}{\textbf{PROP}}
\newcommand{\PRED}{\textbf{PRED}}
\newcommand{\PROPT}{\textbf{PROP2}}
\newcommand{\PREDT}{\textbf{PRED2}}
\newcommand{\PROPo}{\textbf{PROP}\underline{\omega}}
\newcommand{\PREDo}{\textbf{PRED}\underline{\omega}}
\newcommand{\PROPO}{\textbf{PROP}\omega}
\newcommand{\PREDO}{\textbf{PRED}\omega}

\newcommand{\Yes}{\mathtt{Yes}}
\newcommand{\No}{\mathtt{No}}

\theoremstyle{definition}
\setbeamertemplate{theorems}[numbered]
\newtheorem{dfn}{定義}
\newtheorem*{ndfn}{定義}
\newtheorem{axiom}{公理}
\newtheorem{thm}{定理}
\newtheorem{prop}{命題}
\newtheorem{lem}{補題}
\newtheorem{cor}{系}
\newtheorem{ex}{例}
\newtheorem{question}{問}
\newtheorem{caution}{注意}
\newtheorem{notation}{記法}
\newtheorem{conjecture}{予想}


\usepackage{pifont}
\lstset{
    language={C++}, %プログラミング言語によって変える。
    basicstyle={\ttfamily\small},
    %% breaklines=true, %折り返し
}
\lstset{
    basicstyle={\ttfamily},
    identifierstyle={\small},
    commentstyle={\smallitshape},
    keywordstyle={\color{blue}\small\bfseries},
    commentstyle={\color{gray}\small},
    stringstyle={\color{red}\small},
    tabsize=2,
    frame={tb},
    breaklines=true,
    columns=[l]{fullflexible},
    numbers=left,
    xrightmargin=0zw,
    xleftmargin=3zw,
    numberstyle={\scriptsize},
    stepnumber=1,
    numbersep=1zw,
    lineskip=-0.5ex
}

\newenvironment<>{varblock}[2][\textwidth]{%
    \setlength{\textwidth}{#1}
    \begin{actionenv}#3%
        \def\insertblocktitle{#2}%
        \par%
    \usebeamertemplate{block begin}}
    {\par%
        \usebeamertemplate{block end}%
\end{actionenv}}

\newenvironment<>{varalertblock}[2][\textwidth]{%
    \setlength{\textwidth}{#1}
    \begin{actionenv}#3%
        \def\insertblocktitle{#2}%
        \par%
    \usebeamertemplate{block alerted begin}}
    {\par%
        \usebeamertemplate{block alerted end}%
\end{actionenv}}

\newenvironment{scprooftree}[1]%
  {\gdef\scalefactor{#1}\begin{center}\proofSkipAmount \leavevmode}%
  {\scalebox{\scalefactor}{\DisplayProof}\proofSkipAmount \end{center} }

\usetheme{Madrid}
% \usecolortheme{crane}
\usecolortheme{default}
\useinnertheme{}
\useoutertheme{}
\renewcommand{\kanjifamilydefault}{\gtdefault}
\usefonttheme{professionalfonts}
\setbeamertemplate{items}[default]
\setbeamertemplate{navigation symbols}{}
\renewcommand{\baselinestretch}{1}

\setbeamercolor{alerted text}{fg=orange}


\title[内田集合位相ゼミ2]{数理研ゼミ：内田伏一『集合と位相』}
\subtitle{第5章 位相空間}
\author{型推栄}
\institute[M1]{大学院情報理工学研究科 I専攻 岡本吉央研究室 M1}
\date{2026/05/22}

\renewcommand{\thefootnote}{*\arabic{footnote}}
% \setbeamertemplate{enumerate items}{(\arabic{enumi})}

\usetheme{Madrid}
% \usecolortheme{crane}
\useinnertheme{}
\useoutertheme{}
\usefonttheme{professionalfonts}
\setbeamertemplate{items}[default]
\setbeamertemplate{navigation symbols}{}
\renewcommand{\baselinestretch}{1}

\begin{document}
\begin{frame}
    \titlepage
\end{frame}


\begin{frame}[fragile]{本章の目的}
    第4章で距離空間に対して考察した近傍，開集合・閉集合，閉包などの概念を抽象化して\\\alert{位相空間}を導入する

    \bigskip

    \begin{enumerate}
        \item[\S 15] 位相
        \item[\S 16] 近傍系と連続写像
        \item[\S 17] 開基と基本近傍系
        \item[\S 18] 点列連続性
    \end{enumerate}

\end{frame}

\newcommand{\Opens}{\mathscr{O}}
\newcommand{\Closeds}{\mathfrak{A}}
% \renewcommand{\Real}{\mathbf{R}}
% \renewcommand{\Natural}{\mathbf{N}}

\begin{frame}[fragile]{\S 15 位相}
    \begin{block}{位相}
        非空集合$X$の部分集合の族（i.e., $\power X$の部分集合）$\Opens$は，次の条件を満足するとき\\$X$の\alert{位相} (topology) であるという．
       \begin{description}
           \item[$\textbf{O}_1$] $X \in \Opens, \varnothing \in \Opens$.
           \item[$\textbf{O}_2$] $O_1, O_2, \ldots, O_k \in \Opens$ならば$O_1 \cap O_2 \cap \dots \cap O_k \in \Opens$.
           \item[$\textbf{O}_3$] $\qty(O_\lambda \mid \lambda \in \Lambda)$を$\Opens$の元から成る集合系とすれば，$\bigcup_\lambda O_\lambda \in \Opens$.
       \end{description} 
    \end{block}

    \medskip{}

    位相$\Opens$を与えられた集合$X$を\alert{位相空間} (topological space) といい，$(X, \Opens)$で表す．

    各集合$O \in \Opens$を$(X, \Opens)$の\alert{$\Opens$-開集合}ないし単に\alert{開集合} (open set) といい，\\
    $X$の元を$(X, \Opens)$の\alert{点} (point) という．
\end{frame}



\begin{frame}[fragile]{\S 15 位相}
    \begin{block}{}
        非空集合$X$の部分集合の族$\Opens$は，次の条件を満足するとき$X$の\alert{位相}であるという．
       \begin{description}
           \item[$\textbf{O}_1$] $X \in \Opens, \varnothing \in \Opens$.
           \item[$\textbf{O}_2$] $O_1, O_2, \ldots, O_k \in \Opens$ならば$O_1 \cap O_2 \cap \dots \cap O_k \in \Opens$.
           \item[$\textbf{O}_3$] $\qty(O_\lambda \mid \lambda \in \Lambda)$を$\Opens$の元から成る集合系とすれば，$\bigcup_\lambda O_\lambda \in \Opens$.
       \end{description} 
    \end{block}

    \begin{alertblock}{お気持ち}
        位相は\alert{点と点を区別する}もの
        \begin{itemize}
            \item 点$x$と$y$が区別される$\iff$$\exists O \in \Opens,$「$x \in O$ かつ$y \notin O$」または「$x \notin O$かつ$y \in O$」
        \end{itemize}
    \end{alertblock}

    \medskip{}

    $X = \qty{x, y, z},\ \Opens = \qty{X, \qty{x, y}, \qty{z}, \varnothing}$
    \begin{figure}[H]
        \centering
        \begin{tikzpicture}
            \node (x) at (0, 0) {\large $x$};
            \node (y) at (2, 0) {\large $y$};
            \node (z) at (4, 0) {\large $z$};

            \draw (2, 0) circle[x radius=3, y radius=0.8];
            \draw (1, 0) circle[x radius=1.5, y radius=0.5];
            \draw (4, 0) circle[radius=0.5];

            \draw[<->] (x) -- node[above]{\scriptsize 区別不可} (y);
        \end{tikzpicture}
    \end{figure} 
\end{frame}

\begin{frame}[fragile]{\S 15 位相}
    \begin{block}{}
        非空集合$X$の部分集合の族$\Opens$は，次の条件を満足するとき$X$の\alert{位相}であるという．
       \begin{description}
           \item[$\textbf{O}_1$] $X \in \Opens, \varnothing \in \Opens$.
           \item[$\textbf{O}_2$] $O_1, O_2, \ldots, O_k \in \Opens$ならば$O_1 \cap O_2 \cap \dots \cap O_k \in \Opens$.
           \item[$\textbf{O}_3$] $\qty(O_\lambda \mid \lambda \in \Lambda)$を$\Opens$の元から成る集合系とすれば，$\bigcup_\lambda O_\lambda \in \Opens$.
       \end{description} 
    \end{block}

    \medskip{}

    \begin{exampleblock}{例15.1}
        $\Opens = \power X$は$X$の位相である．
        \begin{itemize}
            \item これを\alert{離散位相} (discrete topology) と呼び，$(X, \power X)$を\alert{離散空間}という．
        \end{itemize}
        $\Opens = \qty{X, \varnothing}$も$X$の位相である．
        \begin{itemize}
            \item これを\alert{密着位相} (indiscrete topology) と呼び，$(X, \qty{X, \varnothing})$を\alert{密着空間}という．
        \end{itemize}
    \end{exampleblock}
\end{frame}

\begin{frame}[fragile]{\S 15 位相}
    \begin{block}{}
        非空集合$X$の部分集合の族$\Opens$は，次の条件を満足するとき$X$の\alert{位相}であるという．
       \begin{description}
           \item[$\textbf{O}_1$] $X \in \Opens, \varnothing \in \Opens$.
           \item[$\textbf{O}_2$] $O_1, O_2, \ldots, O_k \in \Opens$ならば$O_1 \cap O_2 \cap \dots O_k \in \Opens$.
           \item[$\textbf{O}_3$] $\qty(O_\lambda \mid \lambda \in \Lambda)$を$\Opens$の元から成る集合系とすれば，$\bigcup_\lambda O_\lambda \in \Opens$.
       \end{description} 
    \end{block}

    \begin{exampleblock}{}
        $\Opens = \power X$を\alert{離散位相}と呼び，$(X, \power X)$を\alert{離散空間}という．
        % $\Opens = \qty{X, \varnothing}$を\alert{密着位相}と呼び，$(X, \qty{X, \varnothing})$を\alert{密着空間}という．
    \end{exampleblock}

    \medskip{}

    $X = \qty{x, y, z},\ \Opens = \qty{X, \qty{x, y}, \qty{y, z}, \qty{x, z}, \qty{x}, \qty{y}, \qty{z}, \varnothing}$
    \begin{figure}[H]
        \centering
        \begin{tikzpicture}
            \node (x) at (0, 0) {\large $x$};
            \node (y) at (2, 0) {\large $y$};
            \node (z) at (1, 1) {\large $z$};

            \draw (x) circle[radius=0.2];
            \draw (y) circle[radius=0.2];
            \draw (z) circle[radius=0.2];

            \draw ($(x)!.5!(y)$) circle[x radius=1.6, y radius=0.4, rotate=0];
            \draw ($(y)!.5!(z)$) circle[x radius=1.4, y radius=0.4, rotate=-45];
            \draw ($(x)!.5!(z)$) circle[x radius=1.4, y radius=0.4, rotate=45];
            \draw (1, 0.4) circle[x radius=2.5, y radius=1.3];

        \end{tikzpicture}
    \end{figure}
\end{frame}

\begin{frame}[fragile]{\S 15 位相}
    \begin{block}{}
        非空集合$X$の部分集合の族$\Opens$は，次の条件を満足するとき$X$の\alert{位相}であるという．
       \begin{description}
           \item[$\textbf{O}_1$] $X \in \Opens, \varnothing \in \Opens$.
           \item[$\textbf{O}_2$] $O_1, O_2, \ldots, O_k \in \Opens$ならば$O_1 \cap O_2 \cap \dots O_k \in \Opens$.
           \item[$\textbf{O}_3$] $\qty(O_\lambda \mid \lambda \in \Lambda)$を$\Opens$の元から成る集合系とすれば，$\bigcup_\lambda O_\lambda \in \Opens$.
       \end{description} 
    \end{block}

    \begin{exampleblock}{}
        % $\Opens = \power X$を\alert{離散位相}と呼び，$(X, \power X)$を\alert{離散空間}という．
        $\Opens = \qty{X, \varnothing}$を\alert{密着位相}と呼び，$(X, \qty{X, \varnothing})$を\alert{密着空間}という．
    \end{exampleblock}

    \medskip{}

    $X = \qty{x, y, z},\ \Opens = \qty{X, \varnothing}$
    \begin{figure}[H]
        \centering
        \begin{tikzpicture}
            \node (x) at (0, 0) {\large $x$};
            \node (y) at (2, 0) {\large $y$};
            \node (z) at (1, 1) {\large $z$};

            % \draw (x) circle[radius=0.2];
            % \draw (y) circle[radius=0.2];
            % \draw (z) circle[radius=0.2];
            %
            % \draw ($(x)!.5!(y)$) circle[x radius=1.6, y radius=0.4, rotate=0];
            % \draw ($(y)!.5!(z)$) circle[x radius=1.4, y radius=0.4, rotate=-45];
            % \draw ($(x)!.5!(z)$) circle[x radius=1.4, y radius=0.4, rotate=45];
            \draw (1, 0.4) circle[x radius=2.5, y radius=1.3];
        \end{tikzpicture}
    \end{figure}
\end{frame}

\begin{frame}[fragile]{}
    \begin{block}{}
        非空集合$X$の部分集合の族$\Opens$は，次の条件を満足するとき$X$の\alert{位相}であるという．
       \begin{description}
           \item[$\textbf{O}_1$] $X \in \Opens, \varnothing \in \Opens$.
           \item[$\textbf{O}_2$] $O_1, O_2, \ldots, O_k \in \Opens$ならば$O_1 \cap O_2 \cap \dots O_k \in \Opens$.
           \item[$\textbf{O}_3$] $\qty(O_\lambda \mid \lambda \in \Lambda)$を$\Opens$の元から成る集合系とすれば，$\bigcup_\lambda O_\lambda \in \Opens$.
       \end{description} 
    \end{block}
    \begin{block}{定理12.2}
        $\Real^n$の開集合系$\Opens = \Opens(\Real^n)$は次の条件を満足する．
       \begin{description}
           \item[(1)] $\Real^n \in \Opens, \varnothing \in \Opens$.
           \item[(2)] $O_1, O_2, \ldots, O_k \in \Opens$ならば$O_1 \cap O_2 \cap \dots O_k \in \Opens$.
           \item[(3)] $\qty(O_\lambda \mid \lambda \in \Lambda)$を$\Opens$の元から成る集合系とすれば，$\bigcup_\lambda O_\lambda \in \Opens$.
       \end{description} 
    \end{block}

    \begin{exampleblock}{例15.2}
        $\Opens(\Real^n)$を$\Real^n$の\alert{通常の位相}と呼ぶ．
    \end{exampleblock} 
\end{frame}

\begin{frame}[fragile]{\S 15 位相}
    \begin{block}{}
        非空集合$X$の部分集合の族$\Opens$は，次の条件を満足するとき$X$の\alert{位相}であるという．
       \begin{description}
           \item[$\textbf{O}_1$] $X \in \Opens, \varnothing \in \Opens$.
           \item[$\textbf{O}_2$] $O_1, O_2, \ldots, O_k \in \Opens$ならば$O_1 \cap O_2 \cap \dots O_k \in \Opens$.
           \item[$\textbf{O}_3$] $\qty(O_\lambda \mid \lambda \in \Lambda)$を$\Opens$の元から成る集合系とすれば，$\bigcup_\lambda O_\lambda \in \Opens$.
       \end{description} 
    \end{block}
    \medskip{}

    距離空間$(X, d)$上の集合$A \subset X$は，$A = A^i$のとき開集合，$A = \overline A$のとき閉集合と呼ぶ．

    \begin{exampleblock}{例15.3}
        距離空間$(X, d)$の開集合系$\Opens$は$X$の位相の一つであり（cf.\ 問13.5），
        これを\\$d$によって定まる\alert{距離位相}と呼ぶ．
        \medskip{}

        また集合$X$上の位相$\Opens$がある距離によって定まる距離位相に一致するとき，\\$\Opens$は\alert{距離化可能} (metrizable) であるという．
    \end{exampleblock}
\end{frame}

\begin{frame}[fragile]{\S 15 位相}
    \begin{block}{}
        非空集合$X$の部分集合の族$\Opens$は，次の条件を満足するとき$X$の\alert{位相}であるという．
       \begin{description}
           \item[$\textbf{O}_1$] $X \in \Opens, \varnothing \in \Opens$.
           \item[$\textbf{O}_2$] $O_1, O_2, \ldots, O_k \in \Opens$ならば$O_1 \cap O_2 \cap \dots O_k \in \Opens$.
           \item[$\textbf{O}_3$] $\qty(O_\lambda \mid \lambda \in \Lambda)$を$\Opens$の元から成る集合系とすれば，$\bigcup_\lambda O_\lambda \in \Opens$.
       \end{description} 
    \end{block}
    \medskip{}

    \begin{exampleblock}{例15.4}
        $(X, \Opens)$: 位相空間， \quad $A \subset X$: 非空
        \medskip{}

        $\Opens_A = \Set{A \cap U | U \in \Opens}$は$A$の位相の一つである．

        これを$A$の上の$\Opens$に関する\alert{相対位相} (relative topology) と呼び，位相空間$(A, \Opens_A)$を\\$(X, \Opens)$の\alert{部分空間} (subspace) という．
    \end{exampleblock}
\end{frame}

\begin{frame}[fragile]{\S 15 位相}
    $(X, \Opens)$: 位相空間

    \medskip{}

    集合$F \subset X$は，$F^c \in \Opens$であるとき$(X, \Opens)$の\alert{閉集合} (closed set) であるという．

    \begin{exampleblock}{問15.3}
        位相空間$(X, \Opens)$の閉集合全体$\Closeds$が次の条件を満足することを確かめよ．
        \begin{description}
            \item[(1)] $X \in \Closeds, \varnothing \in \Closeds$.
            \item[(2)] $F_1, F_2, \ldots, F_k \in \Closeds$ならば$F_1 \cup F_2 \cup \dots F_k \in \Closeds$.
            \item[(3)] $\qty(F_\lambda \mid \lambda \in \Lambda)$を$\Closeds$の元から成る集合系とすれば，$\bigcap_\lambda F_\lambda \in \Closeds$.
        \end{description}
    \end{exampleblock}

    \medskip{}
    \fbox{\large 証明の方針}

        \begin{enumerate}
            \item[(1)] 簡単．
            \item[(2)] ド・モルガンの法則から従う．
            \item[(3)] $\qty(\bigcap_\lambda F_\lambda)^c = \bigcup_\lambda F_\lambda^c$が言えれば十分．
        \end{enumerate}

\end{frame}

\begin{frame}[fragile]{\S 15 位相}
    $(X, \Opens)$: 位相空間， \quad $A \subset X$

    \medskip{}

    $A$に包まれる開集合全体の和集合を$A^i$で表す (i.e., $A^i := \bigcup \Set{O \in \Opens | O \subset A}$) と，\\
    $A^i$自身も$\Opens$に属する．したがって$A^i$は$A$に包まれる最大の開集合である．

    \begin{block}{開核作用子}
        $A^i$を$A$の\alert{内部}または\alert{開核} (interior) といい，$A^i$の点を$A$の\alert{内点} (interior point) という．

        $X$の各部分集合にその内部を対応させる写像
        \begin{equation*} 
            \begin{array}{rccc}
                (-)^i: & \power X & \to & \power X \\
                       & A & \mapsto & A^i
            \end{array}
        \end{equation*}
        を$(X, \Opens)$の\alert{開核作用子} (interior operator) という．
    \end{block}

    \medskip{}

    距離空間のときに考えた内部のイメージでOK

\end{frame}

\begin{frame}[fragile]{\S 15 位相}
    $(X, \Opens)$: 位相空間， \quad $A \subset X$，\quad $A^i := \bigcup \Set{O \in \Opens | O \subset A}$

    \medskip{}

    \begin{block}{定理15.1}
        位相空間$(X, \Opens)$の開核作用子は，次の性質を持つ．
        \begin{description}
            \item[(1)] $X^i = X$.
            \item[(2)] $A^i \subset A$.
            \item[(3)] $(A \cap B)^i = A^i \cap B^i$.
            \item[(4)] $(A^i)^i = A^i$.
        \end{description}
    \end{block}

    \medskip{}

    \fbox{\large 証明} \quad (1), (2), (4) は自明なので (3) のみ示す．

    \begin{description}[$(A \cap B)^i \subset A^i \cap B^i$:]
        \item[$A^i \cap B^i \subset (A \cap B)^i$:] (左辺)\ $\subset A \cap B$で，右辺はそのような開集合の中で最大．
        \item[$(A \cap B)^i \subset A^i \cap B^i$:] (右辺)\ $\subset A^i$,\ (右辺)\ $\subset B^i$から従う．\qed
    \end{description}

\end{frame}

\begin{frame}[fragile]{\S 15 位相}
    $(X, \Opens)$: 位相空間， \quad $A \subset X$

    \medskip{}

    $A$を包む閉集合全体の共通部分を$\bar A$または$A^a$で表す (i.e., $A^a := \bigcap \Set{F \in \Closeds{} | A \subset F}$) と，\\
    $A^a$自身も$\Closeds$に属する．したがって$A^a$は$A$を包む最小の閉集合である．

    \begin{block}{閉包作用子}
        $A^a$を$A$の\alert{閉包} (closure) といい，$A^a$の点を$A$の\alert{触点} (adherent point) という．

        $X$の各部分集合にその閉包を対応させる写像
        \begin{equation*} 
            \begin{array}{rccc}
                (-)^a: & \power X & \to & \power X \\
                       & A & \mapsto & A^a
            \end{array}
        \end{equation*}
        を$(X, \Opens)$の\alert{閉包作用子} (closure operator) という．
    \end{block}

    \medskip{}


\end{frame}

\begin{frame}[fragile]{\S 15 位相}
    $(X, \Opens)$: 位相空間， \quad $A \subset X$，\quad $A^a := \bigcap \Set{F \in \Closeds{} | A \subset F}$

    \medskip{}

    \begin{block}{定理15.2}
        位相空間$(X, \Opens)$の閉包作用子は，次の性質を持つ．
        \begin{description}
            \item[(1)] $\varnothing^i = \varnothing$.
            \item[(2)] $A \subset A^a$.
            \item[(3)] $(A \cup B)^a = A^a \cup B^a$.
            \item[(4)] $(A^a)^a = A^a$.
        \end{description}
    \end{block}

    \medskip{}

    \fbox{\large 証明} \quad 定理15.2の双対．\qed

\end{frame}

\begin{frame}[fragile]{\S 15 位相}
    \begin{exampleblock}{問15.4}
        $(X, d)$を距離空間，$\Opens$を$d$によって定まる距離位相とする．各$A \subset X$について次を示せ
        \begin{enumerate}
            \item \alert{距離空間$(X, d)$における$A$の内部}と，\structure{位相空間$(X, \Opens)$における$A$の内部}は一致する
            \item \alert{距離空間$(X, d)$における$A$の閉包}と，\structure{位相空間$(X, \Opens)$における$A$の閉包}は一致する
        \end{enumerate}
    \end{exampleblock}

    \medskip{}

    \fbox{\large 証明} \quad 1のみ示す．任意の$I \subset A$について「\alert{$A^i = I$}$\iff$\structure{$A^i = I$}」が言えればよい．
   \begin{description}[$\Rightarrow$]
       \item[$(\Rightarrow)$] $I$は\alert{開集合}なので\structure{開集合}でもある．また\structure{開集合}$J \subset A$とその点$x \in J$を任意に取ると，\\
           $x$は$J$の\alert{内点}だから，ある近傍$N(x; \varepsilon) \subset J \subset A$が存在する．
           つまり$x$は$A$の\alert{内点}でもあるので$x \in I$．すなわち$J \subset I$が導かれるから，$I$は$A$に包まれる\structure{開集合}で最大．
       \item[$(\Leftarrow)$] $A$の\alert{内点}$x$を任意に取り，その\alert{近傍}を$J := N(x; \varepsilon) \subset A$とすると，$J$は\structure{開集合}である．\\
           $I$は$A$に包まれる\structure{開集合}で最大だから$J \subset I$．よって$I$は$A$の\alert{内点}を全て含む．\\
           また$I \in \Opens,I \subset A$より$I$の点は全て$A$の\alert{内点}であるから，$I$は$A$の\alert{内点}全体．\qed
   \end{description} 

\end{frame}

\begin{frame}[fragile]{\S 15 位相}
    \begin{exampleblock}{問15.5}
        $(X, \Opens)$を位相空間とする．各$A \subset X$について次を示せ
        \begin{enumerate}
            \item $(A^c)^a = (A^i)^c$
            \item $(A^c)^a = (A^a)^c$
        \end{enumerate}
    \end{exampleblock}

    \medskip{}

    \fbox{\large 証明} \quad 1のみ示す．
   \begin{description}[$(\subset)$]
       \item[$(\subset)$] $((A^c)^a)^c$は開集合である．
       \item[$(\supset)$]
   \end{description} 

\end{frame}

\end{document}
